\section{Proof of Theorem~\ref{thm-adaptive}} \label{app-adaptive}

Fix a PPT adversary $\advA$ making at most $q=q(\secp)$ signing queries, and let $N\le q$ be the number it makes in a given execution. Since $q$ is polynomial and $r$ is super-polynomial, $q\le r$ for all sufficiently large $\secp$. Throughout, $d=\dig(m)$ denotes the encoded digest $\enc{\TLFEv(K,m)}_d$ of a message $m$ under the current pre-hash key $K$.

\heading{From the real game to the lossy chain.} The real game is the EUF-CMA game for $\SchSScheme[\GG,\HF^{\mathsf{sio}}]$, in which the pre-hash key is generated in injective mode. We first add two aborts. The event $\mathsf{Bad}_{\mathsf{rep}}$ occurs if two signing queries use the same slot $u=(d,\salt)$. The event $\mathsf{Bad}_{\mathsf{col}}$ occurs if two distinct messages among the signing queries and the forgery message have equal digests. When either event occurs, the game outputs 0. The games with and without the aborts are identical until an abort occurs, so their output probabilities differ by at most the abort probability in injective mode. In injective mode, distinct messages have distinct digests unless key sampling fails, so $\mathsf{Bad}_{\mathsf{col}}$ requires the failure event, and $\mathsf{Bad}_{\mathsf{rep}}$ requires a repeated salt or the failure event. Hence the two aborts cost at most
$$
\frac{q(q-1)}{2^{\saltlen+1}}+\nu_{\mathsf{inj}}(\secp,r) \;.
$$

Next switch the pre-hash key to lossy mode at the same $r$. The distinguisher for this switch receives $K$, samples $x$ and all PRF keys itself, builds the obfuscated circuit, runs the game including both abort checks, and outputs the game's result. It is PPT, so the switch costs $\AdvTLF{\TLF,r}{\advD_{\mathsf{mode}}}$ for a PPT distinguisher $\advD_{\mathsf{mode}}$. All subsequent games use a lossy key. In particular, in every subsequent game an execution that reaches $\Finalize$ without aborting has all digests of distinct relevant messages distinct, because $\mathsf{Bad}_{\mathsf{col}}$ output 0 otherwise.

\heading{The outer games.} For $0\le i\le q$, Game $\Gm_i$ is given in Figure~\ref{fig-game-adaptive}. It answers the first $i$ signing queries by the trigger procedure and the rest honestly. The salts and honest nonces are sampled before key generation. Moving these independent choices earlier changes no distribution. Every game publishes the original circuit $\CircS$ of Figure~\ref{fig-circuit-salted}. Game $\Gm_0$ is the lossy game with aborts of the previous paragraph, and $\Gm_q$ answers every signing query without the secret exponent.

\begin{figure}[t]
\begin{center}
\fbox{
\begin{minipage}{4.75in}
\begin{tabbing}
123\=123\=123\=\kill
\textbf{Game} $\Gm_i$: \\
\textsc{Initialize}: \\
\> $(K,\mathsl{td})\getsr\TLFKg(\usecp,r,\mathsf{lossy})$ \\
\> For $j\in[q]$: $\salt_j\getsr\bits^{\saltlen} \;;\; r_j\getsr\Zp$ \\
\> $x\getsr\Zp \;;\; X\gets g^x \;;\; k_\ell\getsr\PRFKg_\ell(\usecp)$ for $\ell\in[4]$ \\
\> $\widetilde{\CircS}\getsr\iO(\usecp,\CircS[X,k_1,k_2,k_3,k_4])$ \\
\> $j\gets0 \;;\; U\gets\emptyset \;;\; M\gets\emptyset \;;\; \mathsf{bad}\gets\false$ \\
\> Return $\vkey=(X,K,\widetilde{\CircS})$ \\
\textsc{SignO}$(m)$: \\
\> $j\gets j+1 \;;\; d\gets\dig(m) \;;\; u\gets(d,\salt_j) \;;\; t\gets\enczero{u}$ \\
\> If $u\in U$: $\mathsf{bad}\gets\true$ \\
\> If $\exists\, m'\in M$ with $m'\ne m$ and $\dig(m')=d$: $\mathsf{bad}\gets\true$ \\
\> $U\gets U\cup\{u\} \;;\; M\gets M\cup\{m\}$ \\
\> If $j\le i$: \\
\>\> $c\gets\PRFEv_3(k_3,t) \;;\; s\gets\PRFEv_4(k_4,t) \;;\; R\gets g^sX^{-c}$ \\
\> Else: \\
\>\> $R\gets g^{r_j} \;;\; c\gets\widetilde{\CircS}(R,d,\salt_j) \;;\; s\gets(r_j+cx)\bmod p$ \\
\> Return $(R,s,\salt_j)$ \\
\textsc{Finalize}$(m^*,\sigma^*)$: \\
\> If $\exists\, m'\in M$ with $m'\ne m^*$ and $\dig(m')=\dig(m^*)$: $\mathsf{bad}\gets\true$ \\
\> If $\mathsf{bad}$ or $m^*\in M$: Return $0$ \\
\> Return $\SchVf(\vkey,m^*,\sigma^*)$
\end{tabbing}
\end{minipage}
}
\end{center}
\vspace{-6pt}
\caption{The outer game $\Gm_i$. Every game publishes the original salted hash circuit and uses a lossy pre-hash key; only the signing procedure changes with $i$.}\label{fig-game-adaptive}
\vspace{-3pt}\hrulefill
\vspace{-1ex}
\end{figure}


We next define the temporary circuit used to compare adjacent games. It occurs only in the proof. Fix a slot $u=(d,\salt)$, two distinct group elements $R,W$, and two values $v,c \in \Zp$. Let
$$
t:=\enczero{u},\qquad
T:=\{\encone{R,u},\encone{W,u}\},
\qquad
L:=\{(\encone{R,u},v),(\encone{W,u},c)\} \;.
$$
The two entries of $L$ are sorted by their first coordinates, and $T$ is the corresponding set of first coordinates. The circuit $C_j$ contains $t,L$, keys $k_1\{T\},k_2\{T\}$ punctured at the two strings in $T$, and keys $k_3\{t\},k_4\{t\}$ punctured at $t$.

\begin{figure}[t]
\begin{center}
\fbox{
\begin{minipage}{4.75in}
\begin{tabbing}
123\=123\=123\=\kill
\textbf{Circuit} $C_j[X,t,L,k_1\{T\},k_2\{T\},k_3\{t\},k_4\{t\}](P,d,\salt)$: \\
\> $t' \gets \enczero{d,\salt} \;;\; e'\gets\encone{P,d,\salt}$ \\
\> If $(e',v_0)\in L$ for some $v_0$: Return $v_0$ \comment{two hardwired points} \\
\> If $t'\ne t$: \\
\>\> $c_t \gets \PRFEv_3(k_3\{t\},t') \;;\; s_t \gets \PRFEv_4(k_4\{t\},t')$ \\
\>\> If $P=g^{s_t}X^{-c_t}$: Return $c_t$ \comment{trigger branch} \\
\> $a \gets \PRFEv_1(k_1\{T\},e') \;;\; b \gets \PRFEv_2(k_2\{T\},e')$ \\
\> Return $\big(f(PX^ag^b)+a\big)\bmod p$ \comment{main branch}
\end{tabbing}
\end{minipage}
}
\end{center}
\vspace{-6pt}
\caption{The temporary circuit used when changing the answer to query $j$. It is padded to the same fixed size bound as the other circuits in the proof.}\label{fig-circuit-query}
\vspace{-3pt}\hrulefill
\vspace{-1ex}
\end{figure}


\heading{Changing query $j$.} Fix $j\in[q]$ and restrict to executions with $N\ge j$. We compare $\Gm_{j-1}$ and $\Gm_j$. At each endpoint, first introduce the bad event
$$
\mathsf{Bad}_{\mathsf{eq}} := [R=W],
\qquad
R:=g^{r_j},
\qquad
W:=g^{s_1}X^{-c},
$$
where $t_j=\enczero{d_j,\salt_j}$, $c=\PRFEv_3(k_3,t_j)$, and $s_1=\PRFEv_4(k_4,t_j)$, and output 0 when it occurs. Conditional on the view before query $j$, the point $R$ is uniform and independent of $W$. The two endpoint aborts therefore cost at most $2/p$ in total. Every intermediate game uses the same rule with its current values of $R,c,s_1,W$. In particular, when a PRF challenge changes $c$ or $s_1$, the game recomputes both $W$ and $\mathsf{Bad}_{\mathsf{eq}}$. This change is part of that PRF hybrid. Thus every game in the middle of the chain has $R\ne W$.

The reduction now guesses the digest of query $j$. It computes $S\getsr\TLFEnum(\mathsl{td})$, chooses $\widehat d$ uniformly from the encoded list $\{\enc{y}_d : y \in S\}$ before publishing the verification key, and aborts if $d_j\ne\widehat d$, where $d_j=\dig(m_j)$. Except with probability $\nu_{\mathsf{enum}}(\secp,r)$, the list $S$ has at most $r$ entries and contains every digest the lossy key can produce, so it contains $d_j$. The guess is unused in the two endpoint games and is independent of the adversary's view. Hence for $b\in\{j-1,j\}$,
$$
\Pr[N\ge j \land \Gm_b\Rightarrow1]
\le |S| \cdot \Pr[\widehat d=d_j \land N\ge j \land \Gm_b\Rightarrow1]
+ \nu_{\mathsf{enum}}(\secp,r) \;,
$$
since on successful enumeration $d_j$ occurs in the encoded list and the guess is uniform on it. Thus the guess costs a factor $|S|\le r$, plus one $\nu_{\mathsf{enum}}$ term per endpoint. The salt $\salt_j$ and nonce $r_j$ were sampled, not guessed. This is the step the pre-hash improves. Without it, the guessed value would be the message itself, at a cost of $2^{\msglen}$; with the lossy key, every digest lies in a set of size at most $r$ that the reduction can enumerate.

Set
$$
u=(\widehat d,\salt_j),\qquad
t=\enczero{u},\qquad
R=g^{r_j},\qquad
t_R=\encone{R,u}.
$$
Define the main-branch values
$$
\begin{aligned}
a&=\PRFEv_1(k_1,t_R),&
b&=\PRFEv_2(k_2,t_R),\\
v&=\big(f(RX^ag^b)+a\big)\bmod p,&
s_0&=(r_j+vx)\bmod p,
\end{aligned}
$$
and the trigger values
$$
c=\PRFEv_3(k_3,t),\qquad
s_1=\PRFEv_4(k_4,t),\qquad
W=g^{s_1}X^{-c},\qquad
t_W=\encone{W,u}.
$$
The two candidate records are $(R,v,s_0)$ and $(W,c,s_1)$.

The forward half of the chain is in Table~\ref{tab-adaptive-query}. Each row gives only the change from the preceding game; unmentioned code has the same distribution, including the rule that outputs 0 on the current event $\mathsf{Bad}_{\mathsf{eq}}$. In $\mathsf{H}_8$ we exchange the two triples. We then reverse the preceding changes with $(W,c,s_1)$ as the returned triple.

\begin{table}[t]
\caption{The forward games for query $j$.}\label{tab-adaptive-query}
\vspace{-4pt}
\begin{center}
\small
\begin{tabular}{c@{\hspace{0.12in}}p{2.83in}@{\hspace{0.1in}}p{0.72in}}
\hline
Game & Change from the preceding game & Justification \\
\hline
$\mathsf{H}_0$ & Use $\CircS$ and return $(R,s_0,\salt_j)$ & Start \\
$\mathsf{H}_1$ & Replace $\CircS$ by $C_j$ with true values & iO \\
$\mathsf{H}_2$ & Make the two $\PRF_1$ values on $T$ uniform & $\PRF_1$ \\
$\mathsf{H}_3$ & Make the two $\PRF_2$ values on $T$ uniform & $\PRF_2$ \\
$\mathsf{H}_4$ & Replace $v$ by uniform and discard $a,b$ & Exact \\
$\mathsf{H}_5$ & Sample $s_0,v$ and set $R=g^{s_0}X^{-v}$ & Exact \\
$\mathsf{H}_6$ & Make $\PRF_3(t)$ uniform & $\PRF_3$ \\
$\mathsf{H}_7$ & Make $\PRF_4(t)$ uniform & $\PRF_4$ \\
$\mathsf{H}_8$ & Return $(W,s_1,\salt_j)$ instead of $(R,s_0,\salt_j)$ & Exact \\
\hline
\end{tabular}
\end{center}
\vspace{-5pt}
\end{table}

We justify the changes. First replace the obfuscation of $\CircS$ by an obfuscation of $C_j$ in Figure~\ref{fig-circuit-query}, with $T=\{t_R,t_W\}$ and $L=\{(t_R,v),(t_W,c)\}$. At the true values these circuits compute the same function. On input $(W,\widehat d,\salt_j)$ both return $c$. On input $(R,\widehat d,\salt_j)$ the trigger comparison fails because $R\ne W$, and both return the main-branch value $v$. At every other input whose last two components are $(\widehat d,\salt_j)$, the trigger comparison fails and both use the main branch. At all other inputs, functionality preservation of puncturing applies. The switch is therefore justified by iO, in its auxiliary-state form from Definition~\ref{def-io}. At both the forward and reverse iO switches, the reduction first checks $\mathsf{Bad}_{\mathsf{eq}}$; if it is set, the reduction outputs 0 without invoking iO, or equivalently submits two identical dummy circuits. Thus iO is invoked only on functionally equivalent circuits.

Next replace the outputs of $\PRF_1$ at both strings in $T$ by independent uniform values, and then do the same for $\PRF_2$. Only the values at $t_R$ are used to compute $v$. We puncture at both strings so that the circuit description does not distinguish $R$ from $W$. The puncture set $T$ is sampled independently of $k_1,k_2$, and all other game data is auxiliary state, so Definition~\ref{def-pprf} applies directly.

Now $a,b$ are uniform. For fixed $R$, the element $U:=RX^ag^b$ is uniform and independent of $a$. It follows that $v=f(U)+a$ is uniform and independent of $R$ and of the remaining state. We may replace $v$ by an independent uniform value and discard $a,b$; this change is exact. We then change variables in the first triple: sampling $r_j,v$, setting $R=g^{r_j}$ and $s_0=r_j+vx$, has the same distribution as sampling $s_0,v$ and setting
$$
R=g^{s_0}X^{-v}.
$$
This is also exact.

Next replace $c=\PRFEv_3(k_3,t)$ by a uniform value and then replace $s_1=\PRFEv_4(k_4,t)$ by a uniform value. These changes follow from puncturable-PRF security at $t$. When $c$ or $s_1$ changes, so do $W$, $t_W$, the current event $\mathsf{Bad}_{\mathsf{eq}}$, and the puncture set $T$. The string $t$ is fixed independently of $k_3,k_4$. In the $\PRF_3$ and $\PRF_4$ reductions, the challenge value first determines $W$ and $t_W$; the reduction then generates the independent keys $k_1,k_2$ and punctures them at $T=\{t_R,t_W\}$. The reverse reductions use the same order.

\heading{Exchanging the two triples.} At this point
$$
 (g^{s_0}X^{-v},v,s_0)
\qquad\text{and}\qquad
 (g^{s_1}X^{-c},c,s_1)
$$
are independent and identically distributed, conditioned on their first components being distinct. If $u$ occurred before query $j$, then query $j$ sets $\mathsf{Bad}_{\mathsf{rep}}$; if it occurs again later, that later query sets the same event. Since the final output is then fixed to 0, the game may stop as soon as the event occurs. Thus no other response uses $u$ on an execution that continues.

Before comparing $\mathsf{H}_7$ and $\mathsf{H}_8$, discard the full PRF keys, $r_j,a,b$, and the unused challenge values. Keep only $x$, the punctured keys, $t$, the canonical sets $T,L$, the prior view, and the independent coins for the other signing queries. These values suffice to continue the game: public hash evaluation uses $C_j$; trigger signing at every other pair uses $k_3\{t\},k_4\{t\}$; and honest signing uses $x$ and $C_j$.

The remaining state does not identify either entry of $L$ as the honest or trigger entry, and it contains no separate copy of $s_0$ or $s_1$. The view before query $j$ and the event $d_j=\widehat d$ are unchanged when the two triples are exchanged. Consequently, replacing the response $(R,s_0,\salt_j)$ by $(W,s_1,\salt_j)$ leaves the continuation and the adversary's view unchanged. This step is exact.

We now reverse the preceding changes while returning $(W,s_1,\salt_j)$. Restore the outputs of $\PRF_4$ and $\PRF_3$; reverse the change of variables for the first triple; restore the main-branch value $v$; restore the outputs of $\PRF_2$ and $\PRF_1$; and finally replace $C_j$ by $\CircS$ using iO. Each reverse game samples the hidden values with the same distribution as in the corresponding forward game, and its puncturable-PRF experiment provides the true values. At the endpoint the original circuit is published and query $j$ is answered by the trigger procedure, giving $\Gm_j$.

For each $j$, there are PPT distinguishers $\advD^{\mathsf{io}}_{j,\beta}$ and $\advD^{\mathsf{prf}}_{j,\beta,\ell}$, for $\beta\in\{0,1\}$ and $\ell\in[4]$, such that, on writing
$$
\epsilon_j
:=
\sum_{\beta\in\{0,1\}}
\left(
\AdvIO{}{\advD^{\mathsf{io}}_{j,\beta}}
+\sum_{\ell=1}^4
\AdvPRF{\PRF_\ell}{\advD^{\mathsf{prf}}_{j,\beta,\ell}}
\right),
$$
we have
$$
\left|\Pr[\Gm_{j-1}\Rightarrow1]-\Pr[\Gm_j\Rightarrow1]\right|
\le r\,\epsilon_j+\frac{2}{p}+2\nu_{\mathsf{enum}}(\secp,r) \;.
$$
Summing over the signing queries gives
$$
\left|\Pr[\Gm_0\Rightarrow1]-\Pr[\Gm_q\Rightarrow1]\right|
\le r\sum_{j=1}^q\epsilon_j+\frac{2q}{p}+2q\,\nu_{\mathsf{enum}}(\secp,r) \;.
$$

Both forgery reductions below maintain the set $U$ and emulate the $\mathsf{Bad}_{\mathsf{rep}}$ abort of Figure~\ref{fig-game-adaptive}.

\heading{A main-branch forgery.} We build a CDL adversary $\advB_{\mathsf{main}}$. It receives a challenge $h$, sets $X=h$, generates the four PRF keys and the obfuscated circuit, and answers every signing query by the trigger procedure. Suppose $\advA$ returns a valid forgery $(R^*,s^*,\salt^*)$ on the fresh message $m^*$ and the hash circuit takes the main branch. Set $u^*=(\dig(m^*),\salt^*)$ and compute
$$
a^*=\PRFEv_1(k_1,\encone{R^*,u^*})
\qquad\text{and}\qquad
b^*=\PRFEv_2(k_2,\encone{R^*,u^*}) \;.
$$
The reduction returns
$$
\left(R^*X^{a^*}g^{b^*},\ (s^*+b^*)\bmod p\right).
$$
As in Theorem~\ref{thm-uf}, validity gives the CDL equation. Since $f$ has no zeros, the output is accepted. If $E_{\mathsf{main}}$ is the event that the forgery is valid and takes the main branch, then
$$
\Pr[E_{\mathsf{main}}]
\le \AdvCDL{\GG,g,f}{\advB_{\mathsf{main}}} \;.
$$

\heading{A trigger-branch forgery.} Let $E_{\mathsf{trig}}$ be the event that the forgery is valid and takes the trigger branch. The reduction computes $S\getsr\TLFEnum(\mathsl{td})$ and guesses
$$
\widehat u=(\widehat d,\widehat{\salt})
\getsr \{\enc{y}_d : y \in S\}\times\bits^{\saltlen}
$$
before publishing the verification key, and aborts unless $\widehat u=(\dig(m^*),\salt^*)$. Except with probability $\nu_{\mathsf{enum}}(\secp,r)$, the list $S$ contains $\dig(m^*)$. Since the honest circuit at the endpoint is independent of $\widehat u$, this costs the factor $|S|\cdot2^{\saltlen}\le r\cdot2^{\saltlen}$.

Set $\widehat t=\enczero{\widehat u}$. For this reduction, puncture $k_4$ at $\widehat t$ and use the temporary circuit $C_{\mathsf{forg}}$ in Figure~\ref{fig-circuit-trigger-forgery}. Its final line reuses the main branch of $\CircS$.

\begin{figure}[t]
\begin{center}
\fbox{
\begin{minipage}{4.65in}
\begin{tabbing}
123\=123\=123\=\kill
\textbf{Circuit} $C_{\mathsf{forg}}[X,k_1,k_2,k_3,k_4\{\widehat t\},\widehat t,Z](P,d,\salt)$: \\
\> $t'\gets\enczero{d,\salt} \;;\; c\gets\PRFEv_3(k_3,t')$ \\
\> If $t'=\widehat t$: \\
\>\> If $P=ZX^{-c}$: Return $c$ \comment{guessed pair} \\
\> Else: \\
\>\> $s_t\gets\PRFEv_4(k_4\{\widehat t\},t')$ \\
\>\> If $P=g^{s_t}X^{-c}$: Return $c$ \comment{ordinary trigger branch} \\
\> Run the main branch of $\CircS$ on $(P,d,\salt)$
\end{tabbing}
\end{minipage}
}
\end{center}
\vspace{-6pt}
\caption{The temporary circuit for a trigger-branch forgery. The main branch is copied unchanged from Figure~\ref{fig-circuit-salted}.}\label{fig-circuit-trigger-forgery}
\vspace{-3pt}\hrulefill
\vspace{-1ex}
\end{figure}

If $\widehat s=\PRFEv_4(k_4,\widehat t)$ and $Z=g^{\widehat s}$, then $C_{\mathsf{forg}}$ and $\CircS$ compute the same function, so iO permits the switch. Puncturable-PRF security then replaces $\widehat s$ by a uniform value, so $Z=g^{\widehat s}$ is uniform in $\GG$.

Adversary $\advB_{\mathsf{trig}}$ receives a DL challenge $Y$. It samples $x\getsr\Zp$, sets $X=g^x$, writes $\widehat c:=\PRFEv_3(k_3,\widehat t)$, and builds $C_{\mathsf{forg}}$ with $Z=Y$. This has exactly the distribution of the uniform hybrid because $Y$ is uniform in $\GG$. The reduction answers signing queries with the trigger procedure using the punctured key. It aborts if a signing query has encoding $\widehat t$; when the guess is correct, this requires a signing query whose slot equals $(\dig(m^*),\salt^*)$, and hence a distinct queried message with digest $\dig(m^*)$, which the $\mathsf{Bad}_{\mathsf{col}}$ abort excludes, or the pair $(m^*,\salt^*)$ itself among the queries, which freshness of $m^*$ excludes. If the valid forgery takes the trigger branch at the guessed pair, then
$$
g^{s^*}=(YX^{-\widehat c})X^{\widehat c}=Y,
$$
so $\advB_{\mathsf{trig}}$ returns $s^*$ as the discrete logarithm of $Y$. Thus there are PPT distinguishers $\advD_{\mathsf{trig}}^{\mathsf{io}}$ and $\advD_{\mathsf{trig}}^{\mathsf{prf}}$ such that
$$
\Pr[E_{\mathsf{trig}}]
\le
r\cdot2^{\saltlen}\left(
\AdvDL{\GG,g}{\advB_{\mathsf{trig}}}
+\AdvIO{}{\advD_{\mathsf{trig}}^{\mathsf{io}}}
+\AdvPRF{\PRF_4}{\advD_{\mathsf{trig}}^{\mathsf{prf}}}
\right)+\nu_{\mathsf{enum}}(\secp,r) \;.
$$

\heading{Conclusion.} Combining the abort and mode-switch costs, the changes to the signing responses, and the two forgery cases gives
$$
\begin{aligned}
&\AdvEUFCMA{\SchSScheme[\GG,\HF^{\mathsf{sio}}]}{\advA}\\
&\quad\le
\AdvTLF{\TLF,r}{\advD_{\mathsf{mode}}}
+\AdvCDL{\GG,g,f}{\advB_{\mathsf{main}}}\\
&\qquad+r\cdot2^{\saltlen}\left(
\AdvDL{\GG,g}{\advB_{\mathsf{trig}}}
+\AdvIO{}{\advD_{\mathsf{trig}}^{\mathsf{io}}}
+\AdvPRF{\PRF_4}{\advD_{\mathsf{trig}}^{\mathsf{prf}}}
\right)\\
&\qquad+r\sum_{j=1}^q\epsilon_j
+\frac{q(q-1)}{2^{\saltlen+1}}
+\frac{2q}{p}
+\nu_{\mathsf{inj}}(\secp,r)
+(2q+1)\,\nu_{\mathsf{enum}}(\secp,r) \;.
\end{aligned}
$$
Every reduction in this bound runs in time polynomial in $\secp$ and $r$; the enumeration of the lossy image is the only super-polynomial cost, and the hypothesis of the theorem bounds adversaries of exactly this time. We bound the display under the hypotheses of the theorem. Each CDL, iO, and puncturable-PRF advantage above is at most $\eps$, and so is the DL advantage, by the tight reduction from DL to zero-free CDL given after the theorem statement. Each $\epsilon_j$ contains two iO and eight puncturable-PRF advantages, so $\epsilon_j\le10\eps$. The trivial algorithm that guesses a discrete logarithm has advantage $1/p$, and the same tight reduction gives $1/p\le\eps$. Using $q\le r$, the display is at most
$$
\AdvTLF{\TLF,r}{\advD_{\mathsf{mode}}}
+\left(1+3r\cdot2^{\saltlen}+10r^2+2r\right)\eps
+\frac{q(q-1)}{2^{\saltlen+1}}
+\nu_{\mathsf{inj}}(\secp,r)
+(2q+1)\,\nu_{\mathsf{enum}}(\secp,r) \;.
$$
The first term is negligible because $\TLF$ is secure and $\advD_{\mathsf{mode}}$ is PPT. The second is at most $16\,r^2\,2^{\saltlen}\eps$, which is negligible by hypothesis. The remaining terms are negligible because $\saltlen=\omega(\log\secp)$, $q$ is polynomial, and $\nu_{\mathsf{inj}},\nu_{\mathsf{enum}}$ are negligible at $r$. This proves the theorem. \qed
\adam{Constant check for your pass: each $\epsilon_j$ has $2$ iO and $8$ PRF terms (forward and reverse sandwich); the aggregation constant $16$ above is loose and just needs $r^2 2^{\saltlen}\eps$ absorption; $q\le r$ uses $r$ super-polynomial.}
